\documentclass[11pt]{article}
\usepackage[margin=1in]{geometry}
\usepackage{amsmath,amssymb}
\usepackage{booktabs}
\usepackage{graphicx}
\usepackage{float}
\usepackage{authblk}
\usepackage[numbers,sort&compress]{natbib}
\usepackage{hyperref}
\usepackage{xcolor}
\usepackage{caption}
\usepackage{threeparttable}
\usepackage{multirow}
\usepackage{siunitx}
\sisetup{table-number-alignment=center, detect-weight=true, detect-inline-weight=math}
\captionsetup{font=small,labelfont=bf}
\captionsetup[table]{skip=4pt}
\hypersetup{colorlinks=true,linkcolor=blue!45!black,urlcolor=blue!45!black,citecolor=blue!45!black}

\newcommand{\rhoS}{\rho}

\title{\textbf{Metric-matched ordinal supervision and state-space sequence mixing\\
improve prime-editing efficiency prediction}}

% TODO(authors): replace with the real author list before submission.
\author[1]{Author Name\thanks{Corresponding author.}}
\affil[1]{University of California, Irvine}
\date{}

\begin{document}
\maketitle

% =====================================================================
\begin{abstract}
\noindent
Prime editing installs precise genomic changes without double-strand breaks, but its
efficiency varies by orders of magnitude across pegRNA designs and cellular contexts,
making computational prediction essential for practical design. Current
state-of-the-art predictors embed substantial mechanistic structure derived from the
prime-editing reaction. We asked whether a minimally mechanistic model, given the same
data, could match or exceed them, and if so which modelling choices are responsible.

We present PE-RankFormer, a dual-encoder relational model that couples a
representation of the wild-type/edited sequence pair with a segment-aware pegRNA
representation through bidirectional cross-attention, conditioned on experimental
context. Trained on the identical 297{,}962-row corpus, cross-validation splits and
held-out set used by OptiPrime, PE-RankFormer attains Spearman
$\rho = 0.9079$ on the complete 20{,}509-row held-out benchmark against OptiPrime's
$0.8690$: a difference of $+0.0389$ (95\% CI $[+0.0286,+0.0498]$; ahead in all 5{,}000
protospacer-clustered bootstrap resamples, $p < 0.0002$). Recomputed \emph{within}
each of 14 experimental conditions, where between-study differences in efficiency level
cannot contribute, the margin is larger rather than smaller ($+0.0506$, 95\% CI
$[+0.0345,+0.0670]$, ahead in 14/14 conditions). It is $+0.0804$ on the more
experimentally diverse Kim partition and $+0.0220$ on Liu, but we show that much of
that partition structure follows from a training-weight asymmetry rather than from
architecture: the comparator's own released weights assign the Kim study $2.0\%$ of its
training gradient while the study supplies $55.3\%$ of the held-out rows, and its own
published held-out metric is reported over the Liu conditions, where our margin is
$+0.0220$ ($p=0.028$) rather than $+0.0803$. A
gradient-boosted tree ensemble on engineered design features, trained under the same
protocol and tuned on held-out data, reaches $0.7413$, placing both models well clear of
the standard feature set.

Two components account for the gain, and a factorial analysis separates them. An
\emph{ordinal} outcome head that predicts cumulative threshold indicators rather than
efficiency values is matched to the rank-correlation metric and to the target's heavy
zero inflation, contributing $+0.0092$. A \emph{bidirectional state-space} sequence
mixer replacing self-attention contributes $+0.0192$ --- twice as much. A single model
containing both matches a five-member ensemble. Isotonic calibration fitted purely on out-of-fold predictions
restores an interpretable absolute efficiency estimate while leaving the ranking
unchanged, and on that scale the model also outperforms the baseline: mean absolute
error $0.0478$ against $0.0590$ and RMSE $0.0912$ against $0.1026$, so the advantage is
not confined to the rank metric.

Finally, using 649 technical replicate groups we estimate an empirical
repeatability ceiling for the Kim partition, implying roughly $+0.10$ of genuine
remaining headroom on both the development folds ($+0.1149$, mean of three) and the
held-out set ($+0.1043$), and we report a benchmark caveat: predictors that emit
tied scores in the zero-inflated regime gain $+0.013$--$0.016$ Spearman irrespective
of quality, so comparisons between models with differing tie behaviour are not
automatically like-for-like.
\end{abstract}

\tableofcontents
\newpage

% =====================================================================
\section{Introduction}

Prime editing \citep{anzalone2019} writes user-specified edits into genomic DNA
without double-strand breaks or donor templates. A prime editing guide RNA (pegRNA)
carries a spacer that directs Cas9 nickase to the target, a primer binding site (PBS)
that anneals to the nicked strand, and a reverse-transcription template (RTT) encoding
the desired edit. Subsequent work substantially improved efficiency through nicking
strategies and mismatch-repair modulation \citep{chen2021}, but efficiency still ranges
from undetectable to near-complete conversion depending on the design and the cellular
context. Because experimentally screening pegRNAs is expensive, accurate
\emph{a priori} prediction is a practical bottleneck.

A succession of machine-learning predictors has addressed this. DeepPE and its
successors modelled pegRNA features for specific systems \citep{kim2021}; PRIDICT
introduced sequence-based deep models with product-purity prediction
\citep{mathis2023} and was later extended across chromatin contexts
\citep{mathis2024}; DeepPrime broadened coverage to many editing systems and cell
types \citep{yu2023}. Most recently, OptiPrime aggregated these resources into a large
harmonised corpus and a strongly mechanistic model, establishing the benchmark we adopt
here \citep{hsu2026}.

These models differ in how much prime-editing biochemistry they encode explicitly. The
scientific question we address is whether that mechanistic structure is
\emph{necessary}, or whether a general-purpose relational architecture with sufficient
capacity and the right supervision can recover the same signal from data alone. This
matters beyond benchmark rank: if mechanism can be substituted for by architecture and
objective, prediction can extend to new editing systems without re-deriving reaction
models for each.

\paragraph{Contributions.}
\begin{enumerate}
\item \textbf{A minimally mechanistic model that outperforms a strongly mechanistic
one} under matched data, splits and evaluation: $\rhoS=0.9079$ vs.\ $0.8690$,
$\Delta=+0.0389$, significant under a protospacer-clustered paired bootstrap.
\item \textbf{An ordinal outcome head} matched to the evaluation metric and to the
target's zero inflation, which we show produces \emph{decorrelated errors}, not merely
better ones.
\item \textbf{A bidirectional state-space sequence mixer} which, by a factorial
analysis, is the larger of the two contributions --- correcting the natural assumption
that the novel objective mattered most.
\item \textbf{Monotone calibration} that recovers absolute efficiency at no cost to
rank performance, on which the model also beats the baseline ($19\%$ lower MAE).
\item \textbf{A replicate-based noise ceiling} showing the benchmark is \emph{not}
saturated, together with a tie-structure artifact that affects any Spearman-based
comparison on zero-inflated data.
\end{enumerate}

\paragraph{Scope of the comparison.}
The comparator is a substantially broader piece of work than this one, and we state the
asymmetry plainly rather than let a single number stand for it. OptiPrime predicts
outcomes for nicking-guide (PE3) and dual-pegRNA (twinPE) strategies, exposes
interpretable per-step biochemical pseudorates by construction, and was validated
prospectively in primary human and mouse cells including \emph{in vivo} correction of a
pathogenic allele \citep{hsu2026}. We compare on exactly one of its axes --- ranking
editing efficiency for single-pegRNA designs on its own benchmark --- and we make no
claim on any of the others. Our result is that on \emph{that} axis a model with no
reaction structure is more accurate, which is a statement about how much of the
accuracy of a mechanistic predictor is attributable to its mechanism, not a claim to
supersede it. The mechanistic parameterisation buys generalisation to reaction variants
and interpretability that our model does not offer and does not attempt.

% =====================================================================
\section{Methods}

\begin{figure}[H]
\centering
\includegraphics[width=0.97\textwidth]{../../results/paper/figures/fig1_architecture.pdf}
\caption{\textbf{PE-RankFormer architecture and supervision.}
\textbf{(A)} Overall model. The target locus is encoded as a single stream of
\emph{paired} (wild-type, edited) tokens, so the edit is represented explicitly rather
than inferred by comparing two sequences; the pegRNA is encoded as a nucleotide stream
with an added functional-segment embedding marking spacer, PBS and RTT. Six edit-encoder
and four pegRNA-encoder blocks mix information along each sequence, two bidirectional
cross-attention blocks align the pegRNA against the target, and experimental context
modulates the pooled representation through FiLM. The ordinal head emits $K-1$
cumulative logits whose mean is the predicted normalised rank; a monotone calibration
map converts that to an absolute efficiency (\S\ref{sec:calib}).
\textbf{(B)} Encoder block. The two mixers compared in this work occupy the same slot
with identical residual topology, normalisation placement and FFN width, so a
comparison isolates the mixing mechanism. The S4D kernel mixes by \emph{distance};
self-attention mixes by \emph{content}.
\textbf{(C)} The $K-1$ ordinal thresholds $t_k$, placed at quantiles of the training
efficiency distribution (shown; log counts, clipped at $y=0.6$). The distribution is
heavily right-skewed with a large point mass at zero, which is what motivates
supervising the cumulative distribution rather than the value.
\textbf{(D)} Schematic of the resulting predictions: each design produces a decreasing
curve over thresholds, and its score is the curve's mean.
\textbf{(E)} Source composition. The Schwank-derived study supplies 58.4\% of training
rows and none of the evaluation set, while the Kim-derived study dominates evaluation;
\S\ref{sec:domain} shows that removing the unevaluated study is nonetheless harmful.}
\label{fig:arch}
\end{figure}


\subsection{Data}

We use the training mix underlying OptiPrime, obtained from its authors: 297{,}962
training rows and a 20{,}509-row held-out set (318{,}471 total). Each row records a
pegRNA design, its genomic target, an experimental context and a measured editing
efficiency (fraction of alleles carrying the intended edit), with indel rate where
available. Three source studies contribute (Table~\ref{tab:corpus}). Their proportions
differ markedly between training and evaluation, which becomes analytically important
in \S\ref{sec:domain}.

\begin{table}[H]
\centering
\begin{threeparttable}
\caption{Composition of the benchmark corpus.}
\label{tab:corpus}
\begin{tabular}{l S[table-format=3.1] S[table-format=3.1] S[table-format=2.1] S[table-format=6.0]}
\toprule
& \multicolumn{2}{c}{Share of rows (\%)} & {Zero-mass} & {Rows} \\
\cmidrule(lr){2-3}
Source study & {Training} & {Held-out} & {(\%)} & {(training)} \\
\midrule
Schwank-derived (\texttt{pridict\_pridict2}) & 58.4 & 0.0 & 10.3 & 174067 \\
Liu-derived (\texttt{hsu2026})               & 22.0 & 44.7 &  0.8 & 65594 \\
Kim-derived (\texttt{deepprime})             & 19.6 & 55.3 & 49.9 & 58301 \\
\midrule
Total                                        & 100.0 & 100.0 & 16.0 & 297962 \\
\bottomrule
\end{tabular}
\begin{tablenotes}[flushleft]\footnotesize
\item Held-out set: 20{,}509 rows. ``Zero-mass'' is the percentage of training rows
with editing efficiency exactly zero.
\item The Schwank-derived study supplies the majority of training rows and none of the
evaluation set. Section~\ref{sec:domain} shows that removing it is nonetheless harmful
($-0.0294$ Spearman), so its distributional mismatch is outweighed by the signal its
volume provides.
\end{tablenotes}
\end{threeparttable}
\end{table}

\subsubsection{Sequence representation}

Two token streams are constructed per row.

\textbf{Edit stream.} The wild-type and edited sequences are aligned and encoded as a
single stream of \emph{paired} tokens, each denoting the (wild-type, edited) base pair
at that position. Positions where the two agree therefore receive a different token
from positions where they differ, so the edit is represented explicitly by the pairing
rather than requiring the model to infer it by comparing two separate streams. Maximum
length 100 tokens plus two sentinels.

\textbf{pegRNA stream.} Spacer, PBS and RTT are concatenated into a nucleotide stream
of maximum length 90, accompanied by a parallel \emph{segment} identifier marking which
functional element each position belongs to. Segment identity is supplied as a learned
embedding added to the nucleotide embedding, so the model need not infer element
boundaries.

\textbf{Context.} Cell type, PE system, Cas9 variant, PAM, scaffold, 3$'$ motif and
source study are encoded as categorical embeddings and concatenated into a context
vector.

\subsubsection{Splits and evaluation surfaces}

We adopt OptiPrime's own five-fold cross-validation partition, which is
protospacer-disjoint, and its held-out fold. Because protospacers recur across many
designs, all statistics are computed with the protospacer as the unit of dependence.

Model development used a four-level hierarchy designed to keep the held-out set
uncontaminated:

\begin{enumerate}
\item \textbf{Matched development folds} --- three folds of $\approx$35.7k rows,
constructed to match the held-out set's Liu/Kim composition. Used freely.
\item \textbf{Internal lockbox} --- 17{,}975 rows over 367 protospacers sharing no
protospacer with any development validation set. Consulted once, to gate a shortlist.
\item \textbf{Official five-fold out-of-fold (OOF) predictions} --- each row scored only
by the checkpoint that held its fold out.
\item \textbf{Held-out set} --- consulted four times in total, once per model
generation, each access gated behind an explicit flag and appended to an audit log.
The generation reported here was frozen in writing before its single evaluation, with
the expected result recorded beforehand.
\end{enumerate}

Every development number reported is out-of-fold. The ensemble-selection protocol
(candidate ranking, combination rule, tie-breaking) was committed to version control
before the search was executed.

We state the held-out accounting explicitly because it bears on how the margin should
be read. Table~\ref{tab:main} reports three model generations on the held-out set, so
the set was necessarily consulted at each: rounds~1--4 correspond to four accesses,
recorded in the repository's audit log. Development of each generation therefore
followed a look at the previous generation's held-out score, which is a form of
adaptive selection across generations even though no single model was selected on
held-out data. Two facts bound the exposure. Every architectural and objective decision
within a generation was made on development folds and the once-consulted lockbox, never
on held-out rows; and the comparison against the baseline is a \emph{paired} difference
on identical rows, so a shared optimistic bias in the absolute level does not transfer
to the margin. The post-hoc member analyses in \S\ref{sec:ens-honest} are the one place
we examine held-out structure directly, and they are labelled as such.

\subsection{Resolution of the development evaluation}
\label{sec:resolution}

Effects are reported against a stated resolution rather than a significance test,
because the quantity that matters operationally is whether an effect survives a
re-run. Two calibrations fix it. An unmodified re-run of the same model at a different
batch size moved out-of-fold Spearman by $+0.0018$, and in one round fourteen distinct
candidates all landed within $+0.0017$ to $+0.0020$ of their control. Separately, three
interventions on orthogonal axes each reached $+0.003$ to $+0.004$ on one development
fold and all three changed sign on a second. We therefore treat
$\approx\!0.005$ as the resolution of the development evaluation: effects below it are
reported as directions, and no claim in this paper rests on one. Promotion required
same-sign replication across three development folds \emph{and} a mean of at least
$+0.005$.

\subsection{Model architecture}

PE-RankFormer comprises five components.

\paragraph{(1) Edit encoder.} Six pre-layer-norm blocks over the paired edit stream.

\paragraph{(2) pegRNA encoder.} Four blocks over the segment-aware pegRNA stream.

\paragraph{(3) Bidirectional cross-attention.} Two blocks in which each stream attends
to the other, so pegRNA positions can align to target positions and vice versa. This
is where PBS/RTT--target complementarity is available to the model.

\paragraph{(4) Context conditioning.} Feature-wise linear modulation
\citep[FiLM,][]{perez2018} applied to the pooled representation:
$h' = (1+\gamma(c))\odot h + \beta(c)$, with $\gamma,\beta$ produced by a small MLP on
the context vector.

\paragraph{(5) Output head.} Attention pooling followed by a two-layer MLP; the final
projection depends on the outcome head (\S\ref{sec:heads}).

Base configuration: $d_{\mathrm{model}}=384$, 6 heads, FFN width 1536, dropout $0.10$,
$\approx$25M parameters.

\subsubsection{Sequence mixing: attention versus state space}
\label{sec:mixer}

The blocks in (1) and (2) mix information along the sequence. We compare two mixers
under otherwise identical block topology --- same residual structure, same pre-LN
placement, same FFN width --- so that any difference isolates the mixing mechanism.

\textbf{Self-attention.} Standard multi-head self-attention \citep{vaswani2017}.

\textbf{Bidirectional diagonal state space (S4D).} A diagonal state-space model
\citep{gu2022s4,gu2022s4d} whose convolution kernel is available in closed form: with
diagonal state matrix $A$, step size $\Delta$ and output matrix $C$,
\begin{equation}
K[\ell] \;=\; 2\,\mathrm{Re}\!\left(\sum_n C_n\,
\frac{e^{\Delta A_n}-1}{A_n}\,\bigl(e^{\Delta A_n}\bigr)^{\ell}\right),
\label{eq:s4d}
\end{equation}
applied by FFT convolution. Two independent kernels (forward and time-reversed) are
summed, since editing efficiency depends on context on both sides of the edit. Padded
positions are zeroed before convolution; because the convolution is linear, zeros
contribute nothing to any output, so padding cannot leak into real positions.

The inductive biases differ substantively. Attention performs content-based lookup;
the state-space kernel performs smooth, distance-parameterised mixing that is
translation-equivariant along the sequence.

\subsubsection{Outcome heads}
\label{sec:heads}

\paragraph{Scalar head (baseline).} One pre-sigmoid score trained with Huber loss.

\paragraph{Simplex head.} Every edited locus ends in exactly one of \{unedited,
correctly edited, indel\}, and the assay measures these as proportions. The simplex
head predicts a three-way distribution, respecting that mutual exclusivity by
construction and supervising on the indel signal that a scalar head discards:
\begin{equation}
\mathcal{L}_{\text{simplex}} = -\sum_{k\in\{U,E,I\}} y_k \log p_k ,
\qquad p=\mathrm{softmax}(z)\in\Delta^2 .
\end{equation}
Indel rate is unmeasured for 42.5\% of rows. Rather than imputing zero --- which would
assert a measurement never made --- the loss \emph{marginalises} the indel class for
those rows, renormalising over the two observed outcomes. The ranking score is the
log-odds of a correct edit against no edit, $z_E - z_U$, which is invariant to the
indel logit.

\paragraph{Ordinal head (this work).} The evaluation metric is Spearman correlation,
which depends only on order. We therefore supervise the target's \emph{cumulative
distribution} directly. Fix $K-1$ thresholds $t_k$ at quantiles of the \emph{training}
efficiency distribution and predict one logit per threshold, trained as $K-1$ binary
problems against the cumulative indicators:
\begin{equation}
\mathcal{L}_{\text{ord}}
= \frac{1}{K-1}\sum_{k=1}^{K-1}
\mathrm{BCE}\bigl(z_k,\; \mathbf{1}[y > t_k]\bigr),
\qquad
s(x) = \frac{1}{K-1}\sum_{k=1}^{K-1}\sigma(z_k),
\label{eq:ordinal}
\end{equation}
where $s(x)$ estimates the row's normalised rank. The construction follows
rank-consistent ordinal regression \citep{cao2020}.

Two properties motivate it. It is \emph{metric-matched}: a head estimating the target's
quantile optimises the evaluated quantity, whereas regression or proportion heads are
only incidentally monotone. And it is \emph{robust to the target's skew}: efficiency is
heavily right-skewed, with $16.0\%$ of \emph{training} rows exactly zero overall and
$49.9\%$ within the Kim-derived partition (Table~\ref{tab:corpus}), so a Huber or
proportion objective concentrates its gradient on a few high-efficiency rows, while
$K-1$ balanced binary decisions spread supervision across the distribution. The zero
mass is higher on the surfaces the model is \emph{scored} on --- $26.8\%$ of held-out
rows and $28.4\%$ of matched development-fold rows, both of which exclude the
low-zero-mass Schwank study --- which is why the tie structure of the metric matters
separately (\S\ref{sec:ties}).

Thresholds are computed from training rows only and stored in the checkpoint, so
evaluation reproduces the score mapping exactly. Duplicate quantiles are removed to
keep thresholds strictly increasing; $K=20$ yields 17--18 distinct thresholds in
practice.

Note the ordinal head \emph{discards} the indel signal the simplex head uses. This is
deliberate: it constitutes a different \emph{view} of the same rows, and
\S\ref{sec:ensemble} shows different views are what an ensemble rewards.

\subsection{Training}

AdamW, learning rate $3\times10^{-4}$, weight decay $0.01$, 5\% linear warmup followed
by cosine decay, gradient clipping at $1.0$, batch size 512, bf16 mixed precision, 30
epochs. Checkpoint selection is by validation Spearman.

Early stopping patience is set to the full 30 epochs deliberately. Across all 64
recorded official-fold runs given the full budget, the best epoch had median $24.5$ and
maximum $29$, with $50\%$ peaking in the final five epochs: validation Spearman on this
problem genuinely improves after long plateaus. Simulating a patience-5 rule on the same
runs would have cost a mean of $0.0001$ and at most $0.0031$, so the budget is cheap
insurance rather than a tuned choice.

Batches are formed by a grouped sampler that keeps rows sharing a ranking group
together, which is required for the optional within-group pairwise ranking loss.

\begin{table}[H]
\centering
\begin{threeparttable}
\caption{Architecture and optimisation settings.}
\label{tab:hyper}
\begin{tabular}{l l l l}
\toprule
\multicolumn{2}{l}{\emph{Architecture}} & \multicolumn{2}{l}{\emph{Optimisation}} \\
\midrule
Model width $d_{\mathrm{model}}$ & 384            & Optimiser            & AdamW \\
Attention heads                  & 6              & Learning rate        & $3\times10^{-4}$ \\
FFN width                        & 1536           & Weight decay         & 0.01 \\
Edit-encoder blocks              & 6              & Schedule             & 5\% warmup, cosine \\
pegRNA-encoder blocks            & 4              & Gradient clipping    & 1.0 \\
Cross-attention blocks           & 2              & Batch size           & 512 \\
Dropout                          & 0.10           & Precision            & bf16 \\
Context embedding                & 32             & Epochs               & 30 \\
S4D state dimension              & 64             & Early stopping       & none (see note) \\
Ordinal thresholds $K-1$         & 17--18         & Selection            & best validation $\rhoS$ \\
Parameters                       & $\approx$25\,M & Runtime per model    & $\approx$2\,h, 1 GPU \\
\bottomrule
\end{tabular}
\begin{tablenotes}[flushleft]\footnotesize
\item Patience equals the epoch budget, i.e.\ no early stopping. Across all 64
recorded official-fold runs given the full budget, the best epoch had median $24.5$ and
maximum $29$, with $50\%$ peaking in the final five epochs. Simulating patience-5
selection on the same runs costs a mean of $0.0001$ and at most $0.0031$.
\item Thresholds are quantiles of the training efficiency distribution, de-duplicated
to remain strictly increasing. $K{=}20$ yields 17--18 distinct values because the
target has a large point mass at zero.
\end{tablenotes}
\end{threeparttable}
\end{table}

\subsection{Evaluation and statistics}

The primary metric is Spearman $\rho$ between predicted and observed efficiency,
reported for the full held-out set and separately for the Liu and Kim partitions, since
they differ substantially in composition.

All model comparisons use a \textbf{paired protospacer-clustered bootstrap} with 5{,}000
resamples. Protospacers, not rows, are resampled: many designs share a protospacer, so
a row-level bootstrap would treat correlated observations as independent and materially
understate the intervals. Both models are evaluated on identical rows in each resample.
We report the observed difference, the 95\% percentile interval, the fraction of
resamples in which our model leads, and a two-sided empirical $p$-value.

Because the two source studies differ in efficiency level, pooled rank correlation
benefits from between-study variance that a within-experiment design task would not
supply. We therefore also report every comparison \emph{stratified}: Spearman is
computed inside each (source, cell line, PE system) condition and averaged with weights
equal to condition size, and the clustered bootstrap is applied to that stratified
difference. Strata with fewer than 50 rows or fewer than five distinct target values are
excluded.

\subsection{Baselines}
\label{sec:baselines}

\textbf{OptiPrime.} Run from the released implementation and its five released fold
checkpoints, ensembled as its own prediction script does, on the held-out rows the
authors supplied. The corpus's per-row \texttt{weight} column is a training weight in
that implementation and does not enter its validation metric; \S\ref{sec:domain-weight}
analyses what it implies for the comparison.

\textbf{Gradient-boosted trees on engineered features.} A feature-based reference under
the identical five-fold protocol (\S\ref{sec:gbm}), using 17 engineered design features
--- PBS and RTT length and GC content, four melting temperatures, ViennaRNA folding
energies for the protospacer, RTT, PBS and 3$'$ extension \citep{lorenz2011}, edit type,
edit length, offset from the nick, mismatch count, and a RuleSet3 on-target score
\citep{doench2016} --- concatenated with the same seven categorical context fields the
network receives. Trained with scikit-learn's histogram gradient-boosting regressor.
Four hyperparameter settings were crossed with two target parameterisations (raw
efficiency, and the within-training-fold normalised rank), and the winner was chosen on
held-out Spearman, deliberately giving the baseline an advantage no other model here
receives.

\textbf{Not reported.} PRIDICT2.0 and DeepPrime-FT were investigated and not completed.
Both require reproducing a hand-engineered feature pipeline exactly, including a
separate pretrained on-target cutting model \citep{kim2019}; a subtly mismatched
reimplementation would yield misleading numbers rather than no numbers, and the tree
baseline above covers the same class of information.

\subsection{Ensembling}
\label{sec:ensemble}

Members are combined by \textbf{equal-weight rank averaging}. Members are not mutually
calibrated --- ordinal members emit rank estimates, simplex members emit proportions ---
so averaging raw scores would largely measure scale mismatch. Fitted weights were
evaluated and not adopted: they were unstable across folds and their advantage
($0.9156$ vs.\ $0.9149$ on development folds) was below the resolution of our
evaluation.

Candidate members were assessed by \emph{incremental ensemble gain} under the frozen
combination rule, together with prediction-rank correlation and residual correlation
against the existing ensemble. Residuals are compared in rank space for the same
calibration reason.

\paragraph{Composition of the frozen system.} The final system averages five members:
ordinal$+$S4D, simplex$+$S4D, ordinal with layerwise FiLM, ordinal with a
feature branch, and simplex with layerwise FiLM. Each member is five checkpoints, one
per official fold.
Individual performance and marginal value are given in Table~\ref{tab:members}.

\paragraph{Selection protocol, fixed in advance.} With eight candidate members there
are 255 non-empty subsets, so choosing the best on repeatedly-used development folds
invites a selection bias of the same order as the effects being measured. The protocol
was therefore committed to version control \emph{before} the search was executed:
rank subsets on the development folds; require a candidate to beat the incumbent on
\emph{every} fold rather than on the mean; gate the shortlist on the lockbox, scored
out-of-fold; and on disagreement prefer the lockbox winner, breaking near-ties (within
$0.001$) toward fewer members.

The development winner (four members) and the lockbox winner (five members) disagreed
by $0.001011$ --- eleven millionths past the pre-set near-tie threshold --- so the rule
selected the five-member subset. The two candidates are statistically
indistinguishable, and we followed the rule as written rather than re-arguing it after
seeing the numbers. The lockbox's own unconstrained best subset ($0.9155$) was
\emph{not} adopted, since selecting on the lockbox would convert the only unworn
evaluation surface into a fitted one.

\subsection{Calibration}
\label{sec:calib}

Rank averaging destroys the output scale. We restore it with a monotone map
$g:\text{rank}\mapsto\text{efficiency}$ fitted by isotonic regression
\citep{zadrozny2002} on out-of-fold development predictions only, then applied once to
the frozen held-out predictions. Because $g$ is non-decreasing it cannot invert any
pair, so Spearman is preserved; the map is verified non-decreasing numerically rather
than assumed. This is a change of units on an already-frozen prediction vector, not a
model selection, and no held-out row participates in fitting.

\subsection{Replicate-based noise ceiling}
\label{sec:ceiling}

To determine whether remaining error is reducible we estimate the benchmark's
repeatability. If a design is measured twice in the same condition and the measurements
disagree, no model can predict either better than they predict each other. For two
independent noisy measurements of a true value, reliability
$R = \mathrm{corr}(y_1,y_2)$ implies a ceiling of $\sqrt{R}$ against a single
observation.

Identifying genuine replicates requires care. Grouping on (spacer, RTT, PBS, cell line,
PE system, Cas9) yields 21{,}994 apparently replicated keys, but inspection shows these
differ in 3$'$ motif, epegRNA status and linker --- they are epegRNA versus plain-pegRNA
constructs, i.e.\ \emph{different molecules}. Treating them as replicates gives
$R=0.7987$, almost exactly our model's Kim score, which would falsely indicate
saturation. We therefore require all sixteen design and condition covariates
\emph{and} the target site to match, leaving 649 clean groups.

A second correction is needed. Treating an observed zero as noiseless inflates the
ceiling, but 21.4\% of rows measured at exactly zero have a non-zero replicate (median
$0.0015$): a zero is a \emph{censored} observation. We therefore estimate the ceiling
non-parametrically, drawing synthetic replicates from the empirical distribution of
replicate values observed near each row's efficiency (nearest-neighbour matching on
1{,}298 ordered replicate pairs), which inherits zero-discordance, heteroscedasticity
and boundedness from the data.

\subsection{Implementation and reproducibility}

PyTorch 2.x; all experiments on NVIDIA L40, RTX~PRO~6000 and RTX~2080~Ti GPUs. Every
run records its git commit, configuration, dataset SHA-256 and fold assignment.
Held-out evaluation is gated behind an explicit flag and audit-logged. The code base
carries 214 unit tests covering, among other things, that padding cannot influence
outputs at real positions, that the ordinal score is monotone in every threshold logit,
and that all prior checkpoints load unchanged after each architectural addition.

% =====================================================================
\section{Results}

\subsection{Held-out performance against OptiPrime}

Under matched training data, splits and evaluation, PE-RankFormer attains
$\rhoS=0.9079$ on the complete 20{,}509-row held-out set against OptiPrime's $0.8690$
(Table~\ref{tab:main}, Figure~\ref{fig:bench}).

\begin{table}[H]
\centering
\begin{threeparttable}
\caption{Held-out performance and significance. All models were trained on the
identical corpus and cross-validation splits, and evaluated on identical rows.}
\label{tab:main}
\begin{tabular}{l S[table-format=1.4] S[table-format=1.4] S[table-format=1.4]}
\toprule
& \multicolumn{3}{c}{Spearman $\rhoS$} \\
\cmidrule(lr){2-4}
Model & {All} & {Liu} & {Kim} \\
& {($n{=}20{,}509$)} & {($n{=}9{,}175$)} & {($n{=}11{,}334$)} \\
\midrule
Gradient-boosted trees on engineered features & 0.7413 & 0.6347 & 0.5745 \\
OptiPrime (mechanistic baseline) & 0.8690 & 0.8365 & 0.7320 \\
PE-RankFormer, round 1           & 0.8865 & 0.8349 & 0.7751 \\
PE-RankFormer, round 3           & 0.8933 & 0.8462 & 0.7836 \\
\textbf{PE-RankFormer, final}    & \bfseries 0.9079 & \bfseries 0.8585 & \bfseries 0.8124 \\
\midrule
\multicolumn{4}{l}{\emph{Advantage of the final model over the baseline}} \\
$\Delta\rhoS$ vs.\ OptiPrime     & {$+0.0389$} & {$+0.0220$} & {$+0.0804$} \\
\bottomrule
\end{tabular}
\begin{tablenotes}[flushleft]\footnotesize
\item Liu and Kim denote the two source studies contributing to the held-out set; no
Schwank-derived rows appear in it.
\item The tree baseline is described in \S\ref{sec:gbm}. It is trained under the
identical five-fold protocol on 17 engineered design features plus the same seven
context fields the network receives.
\item Significance for the overall comparison is given in Table~\ref{tab:boot}.
\item The absolute figure for the final model has a bootstrap 95\% interval of
$[0.8955, 0.9173]$, with 90.0\% of resamples exceeding $0.90$; the margin over the
baseline is therefore a stronger claim than the absolute level.
\end{tablenotes}
\end{threeparttable}
\end{table}

\begin{figure}[H]
\centering
\includegraphics[width=\textwidth]{../../results/paper/figures/fig2_benchmark.pdf}
\caption{\textbf{Held-out benchmark result.}
\textbf{(A)} Spearman $\rho$ on the complete held-out set and on each source partition,
for a feature-based reference, the mechanistic baseline, and successive generations of
our model. All were trained on the identical corpus and cross-validation splits and
evaluated on identical rows; the dotted line marks the baseline's overall performance.
The ordering --- engineered features, then mechanistic reaction model, then learned
sequence representation --- places the margin in context: the tree baseline is
$0.128$ behind the mechanistic model, which is in turn $0.039$ behind ours.
\textbf{(B)} Distribution of the paired difference in Spearman $\rho$ under a
protospacer-clustered bootstrap (5{,}000 resamples over 750 clusters). Protospacers
rather than rows are resampled because many designs share a target; the solid red line
is the observed difference, dashed lines the 95\% percentile interval, and the black
line marks zero.
\textbf{(C)} The same comparison \emph{within} each experimental condition
($n\geq300$), where between-study differences in efficiency level cannot contribute to
the correlation. Each line joins the two models on one condition. Our model leads in
all 14, and the $n$-weighted margin is larger within condition ($+0.0506$) than pooled
($+0.0389$), so pooling makes the reported difference conservative
(\S\ref{sec:stratified}). Absolute levels are lower for both models here, and lowest in
the conditions with the highest fraction of exactly-zero measurements
(cf.\ Figure~\ref{fig:ceiling}B).}
\label{fig:bench}
\end{figure}

The paired protospacer-clustered bootstrap (Table~\ref{tab:boot},
Figure~\ref{fig:bench}B) gives $\Delta\rhoS = +0.0389$ with a 95\% interval of
$[+0.0286,+0.0498]$; PE-RankFormer leads in every one of 5{,}000 resamples.

\begin{table}[H]
\centering
\begin{threeparttable}
\caption{Paired significance testing on the held-out set.}
\label{tab:boot}
\begin{tabular}{l S[table-format=+1.4] c S[table-format=1.3] l}
\toprule
Comparison & {$\Delta\rhoS$} & {95\% CI} & {Fraction ahead} & {$p$} \\
\midrule
Final vs.\ OptiPrime, all rows   & +0.0389 & $[+0.0286,\,+0.0498]$ & 1.000 & $<0.0002$ \\
\quad --- Kim partition only     & +0.0803 & $[+0.0632,\,+0.0985]$ & 1.000 & $<0.0002$ \\
\quad --- \textbf{Liu partition only} & +0.0220 & $[+0.0027,\,+0.0427]$ & 0.986 & 0.028 \\
Final vs.\ round 3     & +0.0146 & $[+0.0113,\,+0.0184]$ & 1.000 & $<0.0002$ \\
Round 3 vs.\ OptiPrime & +0.0243 & $[+0.0138,\,+0.0346]$ & 1.000 & $<0.0002$ \\
\bottomrule
\end{tabular}
\begin{tablenotes}[flushleft]\footnotesize
\item Paired bootstrap over 5{,}000 resamples of 750 \emph{protospacer clusters}.
Protospacers rather than rows are resampled because many designs share a target;
a row-level bootstrap would treat correlated observations as independent and
understate the intervals.
\item Both models are scored on identical rows within each resample. ``Fraction
ahead'' is the proportion of resamples in which PE-RankFormer leads; $p$ is the
two-sided empirical value, bounded below by $1/5000$.
\item The partition rows resample the protospacer clusters present in that partition
(601 for Kim, 150 for Liu). \textbf{The Liu result is the weakest claim in the paper}
and we give it prominence rather than only the pooled figure: the comparator's own
published held-out metric is reported over the Liu conditions
\citep{hsu2026}, so Liu is the surface it selected, and there our advantage is
$+0.0220$ with a lower bound of $+0.0027$ at $p=0.028$. The pooled and Kim results are
far stronger, but Kim is the partition the comparator trains on at one tenth weight
(\S\ref{sec:domain-weight}), and the full-set evaluation is our extension of its
benchmark rather than its own headline surface.
\end{tablenotes}
\end{threeparttable}
\end{table}



We note explicitly that the absolute figure's own 95\% interval is $[0.8955,0.9173]$,
with 90.0\% of resamples exceeding $0.90$. The \emph{margin} over OptiPrime is a
considerably stronger claim than the absolute level.

\subsection{A feature-based baseline places both models}
\label{sec:gbm}

A single external comparator leaves open the question of how much of the accuracy comes
from sequence modelling at all. We therefore trained a gradient-boosted tree ensemble
under the identical five-fold protocol on 17 engineered design features --- PBS and RTT
lengths and GC content, four melting temperatures, ViennaRNA folding energies
\citep{lorenz2011}, edit type, edit length, offset from the nick, mismatch count and a
RuleSet3 on-target score \citep{doench2016} --- concatenated with the same seven
categorical context fields the network receives. This is the same class of information
the feature-based predictors we could not run end-to-end rely on
(\S\ref{sec:limitations-baselines}), so it stands in for them.

The baseline was given every advantage: four hyperparameter settings crossed with two
target parameterisations (raw efficiency and the metric-matched rank transform), with
the winner selected on \emph{held-out} Spearman --- the only model in this paper allowed
to tune on the test set. Its best configuration reaches $\rhoS = 0.7413$
(Table~\ref{tab:main}), against OptiPrime's $0.8690$ and PE-RankFormer's $0.9079$; within
condition it reaches $0.5460$ against PE-RankFormer's $0.8111$.

Two things follow. The sequence model earns its parameters: $+0.167$ pooled and
$+0.265$ within condition over a tuned tree ensemble on the standard feature set is far
outside anything attributable to capacity or tuning. And the mechanistic comparator is
itself well clear of that feature set ($+0.128$), so its reaction structure is doing
real work relative to hand-engineered summaries --- which is the more informative way to
read our margin over it. The ordering is: engineered features $<$ mechanistic reaction
model $<$ learned sequence representation.

\subsection{The margin is not an artifact of pooling}
\label{sec:stratified}

Pooled Spearman on the held-out set ($0.9079$) exceeds both partition scores (Liu
$0.8585$, Kim $0.8124$). That is possible because the two studies differ in efficiency
level, so between-study variance contributes signal that no within-condition design
task would supply, and it means the pooled figure overstates how well either model
orders designs \emph{within} an experiment. The model is also given source study as one
of its seven context fields, which a reader may reasonably regard as a dataset
identifier rather than a biological covariate.

Both observations bear on the absolute level. Neither weakens the margin, and the
stratified analysis shows why (Table~\ref{tab:stratified}): recomputing the comparison
\emph{inside} each experimental condition, where no between-condition offset can
contribute, gives $+0.0506$ rather than $+0.0389$. The margin is larger, not smaller,
when the pooling advantage is removed, and PE-RankFormer leads in all 14 conditions
with $n\geq300$ and in 100\% of protospacer-clustered bootstrap resamples. Pooling
inflates the absolute level for both models while making the reported difference
\emph{conservative}.

\begin{table}[H]
\centering
\begin{threeparttable}
\caption{The advantage over the baseline under three levels of stratification.}
\label{tab:stratified}
\begin{tabular}{l S[table-format=2.0] S[table-format=1.4] S[table-format=1.4] S[table-format=+1.4] c}
\toprule
Comparison & {Strata} & {OptiPrime} & {PE-RankFormer} & {$\Delta\rhoS$} & {95\% CI} \\
\midrule
Pooled (as reported in Table~\ref{tab:main}) & {1} & 0.8690 & 0.9079 & +0.0389 & $[+0.0286,+0.0498]$ \\
Within source study, $n$-weighted            & {2} & 0.7788 & 0.8331 & +0.0543 & {---} \\
\textbf{Within condition}, $n$-weighted      & {14} & 0.7605 & 0.8111 & \bfseries +0.0506 & $[+0.0345,+0.0670]$ \\
\bottomrule
\end{tabular}
\begin{tablenotes}[flushleft]\footnotesize
\item A condition is a (source study, cell line, PE system) cell. Strata are weighted by
row count; strata with fewer than 50 rows or fewer than 5 distinct target values are
excluded, which removes none of the 14 reported conditions.
\item The interval on the stratified difference is a paired bootstrap over 5{,}000
resamples of the same 750 protospacer clusters used in Table~\ref{tab:boot};
PE-RankFormer leads in 100\% of resamples. Per-condition margins range from $+0.0084$
(Liu, HEK293T, PE2) to $+0.1295$ (Kim, MDA-MB-231, PE2) and are shown in
Figure~\ref{fig:bench}C.
\item The absolute levels fall under stratification for both models because
between-condition differences in efficiency level no longer contribute to the
correlation. This is the quantity relevant to prospective design, where a user ranks
candidate pegRNAs within one intended experiment.
\end{tablenotes}
\end{threeparttable}
\end{table}

\subsection{A training-weight asymmetry accounts for much of the partition margin}
\label{sec:domain-weight}

The margin is largest on the Kim partition ($+0.0804$), and the natural reading --- that
a mechanistic prior generalises less well across cell lines --- turns out not to be the
most parsimonious one. The corpus the comparator's authors supplied carries a per-row
\texttt{weight} column, and the released training code consumes it as a per-row
\emph{loss} weight, forming the objective as
$\sum_i w_i \ell_i / n_{\text{batch}}$.\footnote{\texttt{reaction/rx\_model.py}, in the
released implementation. The same column is passed to the validation loader but is used
there only to identify padding rows, so it does not enter the reported validation
metric.} Every one of the 69{,}635 Kim rows carries weight exactly $0.1$, against a mean
of $0.97$ for Liu and $0.52$ for Schwank, the latter two varying continuously row by row.

The consequence is a sharp mismatch between what the baseline was trained to fit and
what the benchmark scores (Table~\ref{tab:weights}). Under its own weights the
comparator allocates $2.0\%$ of its effective training gradient to the Kim study, which
supplies $55.3\%$ of the held-out rows. We train uniformly, giving Kim $19.6\%$. Any
uniformly-weighted model therefore enjoys an advantage on this benchmark's dominant
partition that owes nothing to architecture or objective.

The weighting is also internally consistent with how the comparator reports itself. Its
published held-out figure is given over the Liu conditions \citep{hsu2026}, not over the
pooled set; treating Kim as low-weight auxiliary training data and evaluating on Liu is
a coherent pair of choices. The pooled 20{,}509-row evaluation is therefore \emph{our}
extension of its benchmark, not its own headline surface, and the two surfaces give very
different pictures of the margin: $+0.0803$ on Kim ($p<0.0002$) against $+0.0220$ on Liu
($p=0.028$, lower bound $+0.0027$; Table~\ref{tab:boot}). Both are reported, and we
regard the Liu figure as the conservative statement of the result.

\begin{table}[H]
\centering
\begin{threeparttable}
\caption{Effective training weight by source study, against held-out composition.}
\label{tab:weights}
\begin{tabular}{l S[table-format=2.1] S[table-format=2.1] S[table-format=2.1] S[table-format=1.2]}
\toprule
& \multicolumn{2}{c}{Share of training gradient (\%)} & {Share of} & {Mean row} \\
\cmidrule(lr){2-3}
Source study & {OptiPrime weights} & {Uniform (ours)} & {held-out (\%)} & {weight} \\
\midrule
Schwank-derived & 73.2 & 58.4 &  0.0 & 0.52 \\
Liu-derived     & 24.8 & 22.0 & 44.7 & 0.97 \\
Kim-derived     &  2.0 & 19.6 & 55.3 & 0.10 \\
\bottomrule
\end{tabular}
\begin{tablenotes}[flushleft]\footnotesize
\item Shares are computed over the 297{,}962 training rows, as
$\sum_{i \in \text{study}} w_i / \sum_i w_i$.
\item Every Kim row carries weight $0.1$ exactly --- a flat per-study factor, not a
per-row precision estimate. Liu and Schwank weights vary continuously (713 and 2{,}685
distinct values), consistent with a depth- or precision-derived quantity clipped to
$[0.1, 1]$.
\item We infer that the released checkpoints were trained with these weights, since the
released code consumes this column from the files the authors supplied. We have not
been able to verify it independently, and flag it as an inference the comparator's
authors could confirm or correct.
\end{tablenotes}
\end{threeparttable}
\end{table}

WEIGHTEXPT

The finding does not withdraw the comparison, and we retain the uniformly-weighted model
as the primary result: training on all available rows at equal weight is the default a
practitioner would adopt, it is the choice our own sweeps support
(\S\ref{sec:domain}, where re-weighting toward the evaluated distribution was harmful
and removing the unevaluated study cost $-0.0294$), and both models are scored on
identical rows either way. What the finding does is relocate part of the explanation.
The pooled margin is a genuine like-for-like result under each model's own training
recipe; the \emph{partition} structure of that margin is substantially a consequence of
a data-curation decision, and should not be read as evidence about mechanism versus
architecture.

\subsection{Which components are responsible}

Because the ordinal head and the state-space mixer were introduced together, we
separate them with a $2\times2$ factorial over out-of-fold development predictions
(Figure~\ref{fig:fact}).

\begin{table}[H]
\centering
\begin{threeparttable}
\caption{Objective $\times$ architecture factorial.}
\label{tab:factorial}
\begin{tabular}{l S[table-format=1.4] S[table-format=1.4] S[table-format=+1.4]}
\toprule
Sequence mixer & {Simplex head} & {Ordinal head} & {Ordinal $-$ simplex} \\
\midrule
Self-attention (Transformer) & 0.8798 & 0.8911 & +0.0114 \\
Bidirectional S4D            & 0.9011 & \bfseries 0.9082 & +0.0071 \\
\midrule
S4D $-$ Transformer & {$+0.0213$} & {$+0.0171$} & \\
\bottomrule
\end{tabular}
\begin{tablenotes}[flushleft]\footnotesize
\item Out-of-fold development Spearman $\rhoS$, averaged over three matched folds.
\item Main effect of \textbf{objective} $=+0.0092$; main effect of
\textbf{architecture} $=+0.0192$. Both exceed the $\approx0.005$ resolution of the
development evaluation (\S\ref{sec:resolution}) and architecture contributes roughly
twice what the objective does. The \textbf{interaction} $=-0.0043$ falls \emph{below}
that resolution and is therefore reported as a direction, not an established effect.
\item These are observational cells rather than a purpose-built factorial: three
derive from one training round and the Transformer/simplex cell is an earlier frozen
baseline. All four share corpus, splits, out-of-fold scoring, capacity and optimiser
recipe. The interaction term is the quantity most exposed to that mismatch, which is
a second reason not to rest any claim on it.
\end{tablenotes}
\end{threeparttable}
\end{table}

\begin{figure}[H]
\centering
\includegraphics[width=\textwidth]{../../results/paper/figures/fig3_factorial.pdf}
\caption{\textbf{Separating the two contributions.}
\textbf{(A)} Out-of-fold development Spearman for each combination of outcome head and
sequence mixer. All four cells share corpus, splits, out-of-fold scoring, capacity and
optimiser recipe.
\textbf{(B)} Main effects and interaction derived from (A). The architecture change
contributes roughly twice what the objective change does, and the negative interaction
indicates the two partly capture the same improvement rather than stacking.
\textbf{(C)} Why the ordinal head is nonetheless valuable: residual correlation with
the existing ensemble, for every candidate model of each head type (points; bars are
group means). Ordinal-head models are systematically more decorrelated than models
that vary seed, capacity or data subsample --- they make \emph{different errors}, which
is what determines a member's value in an ensemble (\S\ref{sec:ensemble}). The group
means differ by about $0.05$; the ranges touch at one point, so the claim is about the
axis of variation rather than about any individual model. Lower is more complementary.}
\label{fig:fact}
\end{figure}

Three observations follow.

\textbf{Architecture is the larger lever.} The state-space mixer contributes $+0.0192$
against the ordinal objective's $+0.0092$. This runs against the intuitive reading ---
the ordinal head is the more novel component and produced the more visible qualitative
change --- and we record the correction explicitly.

\textbf{The two appear to combine sub-additively} ($-0.0043$): ordinal supervision is
worth $+0.0114$ on a Transformer but only $+0.0071$ on a state-space backbone, which
would indicate the mechanisms partly capture the same improvement rather than stacking.
We flag this as suggestive only. The interaction is smaller than the $\approx0.005$
resolution of the development evaluation and it is the cell combination most exposed to
the observational design, so we do not claim sub-additivity as a result; we note that
the data are consistent with it and that nothing in our conclusions depends on it.

\textbf{A single model suffices.} Combining both yields $\rhoS=0.9082$ out-of-fold,
exceeding a three-member ensemble of earlier models ($0.8982$).

These are observational cells rather than a purpose-built factorial: three derive from
one training round and the Transformer/simplex cell is an earlier frozen baseline. They
share corpus, splits, out-of-fold scoring, capacity and optimiser recipe. Both main
effects clear the $\approx0.005$ resolution of the development evaluation
(\S\ref{sec:resolution}); the interaction term does not, and is additionally the
quantity most exposed to the observational design.

\subsection{The ordinal head produces decorrelated errors}
\label{sec:decorrelation}

The hypothesis motivating the ordinal head was that a different loss geometry would
produce different \emph{errors}, not merely better ones. The corresponding diagnostic is
residual correlation with the existing ensemble, measured over all fifteen candidates of
the round-4 search. Across the seven ordinal-head candidates it averages $0.705$
(range $0.686$--$0.735$); across the six candidates that vary seed, capacity or data
subsample instead of the objective it averages $0.756$ (range $0.686$--$0.781$).

The separation is systematic but not disjoint, and we state it that way rather than as a
clean split: one bagged simplex run reaches $0.686$, inside the ordinal range, and its
incremental ensemble gain ($+0.0029$) was correspondingly middling rather than
top-ranked. The transferable claim is the comparison between \emph{axes} of variation,
not a threshold on any single model --- changing the objective moves residual
correlation by about $0.05$, whereas changing seed, width or data subsample leaves it
where it was.

This distinction has practical consequences for ensembling
(Figure~\ref{fig:members}). The best combination pairs models differing in
\emph{both} objective and architecture: the Transformer/simplex and state-space/ordinal
diagonal has the lowest residual correlation of any pair measured ($0.6453$), whereas
two models sharing an architecture correlate at $0.7641$.

\begin{figure}[H]
\centering
\includegraphics[width=\textwidth]{../../results/paper/figures/fig4_ensemble.pdf}
\caption{\textbf{Ensemble structure on held-out data.}
\textbf{(A)} Residual rank correlation between the five members. Members sharing an
architecture are the most redundant pair; members differing in both objective and
architecture are the most complementary.
\textbf{(B)} Each member's individual held-out Spearman. The dashed line marks the
five-member ensemble, which the strongest single member matches.
\textbf{(C)} Change in held-out Spearman when each member is removed. Two members
contribute negatively, indicating the ensemble is at or past its useful size. Panels
(B) and (C) are post-hoc analyses on held-out data and did not inform the frozen model
choice (\S\ref{sec:ens-honest}).}
\label{fig:members}
\end{figure}

\subsection{Decorrelation is necessary but not sufficient}

A natural reading of the previous section is that diversity should be maximised. It
should not. We tested a model trained with an auxiliary pairwise ranking loss that
proved to be the \emph{most decorrelated} candidate we measured --- lowest prediction
correlation and lowest residual correlation of any model in the study --- and it
\emph{degraded} the ensemble ($-0.0037$), because at $\rhoS=0.8244$ standalone it could
not carry its weight in an equal-weight average. A quantile-regression head reproduced
the same pattern on an unrelated objective: lowest residual correlation
($0.6645$), lowest standalone score ($0.8840$), lowest incremental gain.

The operative criterion is therefore that a member must be \textbf{both decorrelated
and competent}. Decorrelation distinguishes useful from redundant \emph{among models of
comparable strength}; it cannot rescue a weak model.

\subsection{Ensemble versus single model}
\label{sec:ens-honest}

\begin{table}[H]
\centering
\begin{threeparttable}
\caption{Ensemble members: individual performance and marginal value on the held-out set.}
\label{tab:members}
\begin{tabular}{l S[table-format=1.4] S[table-format=1.4] S[table-format=1.4] S[table-format=+1.4]}
\toprule
& \multicolumn{3}{c}{Individual $\rhoS$} & {Marginal} \\
\cmidrule(lr){2-4}
Member & {All} & {Liu} & {Kim} & {contribution} \\
\midrule
ordinal $+$ S4D       & \bfseries 0.9082 & 0.8576 & \bfseries 0.8151 & +0.0031 \\
simplex $+$ S4D       & 0.9017 & 0.8563 & 0.7982 & +0.0013 \\
ordinal $+$ layerwise & 0.9005 & 0.8394 & 0.8082 & +0.0006 \\
ordinal $+$ features  & 0.8947 & 0.8261 & 0.8027 & -0.0007 \\
simplex $+$ layerwise & 0.8888 & 0.8412 & 0.7787 & -0.0013 \\
\midrule
\textbf{Ensemble} (equal-weight rank average) & \bfseries 0.9079 & \bfseries 0.8585 & 0.8124 & \\
\bottomrule
\end{tabular}
\begin{tablenotes}[flushleft]\footnotesize
\item Each member is five checkpoints, one per official fold; every row is scored only
by the checkpoint that held its fold out.
\item ``Marginal contribution'' is the change in ensemble Spearman when that member is
removed, so a positive value means the ensemble is worse without it.
\item The strongest single member matches the full ensemble and two members contribute
negatively. These are post-hoc analyses on held-out data and did not inform the frozen
model choice.
\item The ordinal$+$S4D member scores $0.9082$ here and also $0.9082$ in
Table~\ref{tab:factorial}, which reports out-of-fold \emph{development} Spearman. This
is a coincidence of rounding, not a transcription error: the two are computed from
different files over disjoint row sets ($20{,}509$ held-out rows versus $107{,}041$
development rows, intersection empty) and differ at the fifth decimal place,
$0.908155$ against $0.908179$. We flag it because the agreement to four decimals
otherwise invites the reasonable suspicion that one number has been copied into two
places.
\end{tablenotes}
\end{threeparttable}
\end{table}

The frozen system is a five-member rank-average ensemble. Post-hoc analysis on the
held-out set shows its strongest single member attains $0.9082$ against the ensemble's
$0.9079$, and that two members contribute negatively (Table~\ref{tab:members}; Figure~\ref{fig:members}C).
On development folds the ensemble was worth $+0.0074$ over its best member; on held-out
that advantage is $-0.0003$.

We report this rather than presenting the ensemble as superior. All differences
involved ($0.0003$--$0.0013$) are far below the $\approx0.005$ resolution of our
development evaluation, so the defensible statement is that the ensemble and its best
single member are \emph{indistinguishable} on this benchmark. For practical deployment
the single ordinal-SSM model is preferable: equal accuracy from one model rather than
twenty-five checkpoints. We did not re-select on held-out data to make that
substitution official, because a model chosen by inspecting the test set carries none of
the guarantees the frozen one does.

\subsection{Calibration restores absolute efficiency, and it beats the baseline there too}
\label{sec:calib-results}

Rank averaging yields excellent orderings but no interpretable scale. Isotonic
calibration fitted on out-of-fold development predictions (\S\ref{sec:calib}) restores
one, leaving Spearman unchanged at $0.9079$ (Table~\ref{tab:calib},
Figure~\ref{fig:calib}).

The comparison that establishes whether the recovered scale is any \emph{good} is
against the baseline on identical rows, not against our own uncalibrated output: a rank
average is not measured in efficiency units, so its nominal MAE of $0.3869$ quantifies
a change of units rather than an error. Against OptiPrime the calibrated model is ahead
on absolute accuracy as well as on ranking --- MAE $0.0478$ against $0.0590$
($-19\%$), RMSE $0.0912$ against $0.1026$ ($-11\%$), Pearson $0.8637$ against $0.8270$
--- so the advantage is not confined to the rank metric the model was designed around.

\begin{table}[H]
\centering
\begin{threeparttable}
\caption{Absolute and rank accuracy on the held-out set. Monotone calibration recovers
an efficiency scale, on which the model also outperforms the baseline.}
\label{tab:calib}
\begin{tabular}{l S[table-format=1.4] S[table-format=1.4] S[table-format=1.4] S[table-format=1.4]}
\toprule
Output & {Spearman $\rhoS$} & {Pearson $r$} & {MAE} & {RMSE} \\
\midrule
OptiPrime (mechanistic baseline)        & 0.8690 & 0.8270 & 0.0590 & 0.1026 \\
\midrule
PE-RankFormer, rank average             & 0.9079 & 0.7278 & \multicolumn{2}{c}{\emph{not in efficiency units}} \\
\textbf{PE-RankFormer $+$ isotonic}     & \bfseries 0.9079 & \bfseries 0.8637 & \bfseries 0.0478 & \bfseries 0.0912 \\
\midrule
Advantage over baseline & {$+0.0389$} & {$+0.0367$} & {$-0.0112$} & {$-0.0114$} \\
\bottomrule
\end{tabular}
\begin{tablenotes}[flushleft]\footnotesize
\item All three rows are scored on the identical 20{,}509 held-out rows.
\item The map is fitted by isotonic regression on out-of-fold development predictions
only and applied once to the frozen held-out predictions; no held-out row participates
in fitting.
\item Because the map is non-decreasing it inverts no pair, so Spearman is preserved.
The residual change of $2\times10^{-5}$ arises because isotonic regression is
piecewise constant and therefore creates ties.
\item The uncalibrated rank average lies on an arbitrary $[0,1]$ scale, so its MAE and
RMSE would measure the choice of units rather than an error; we omit them rather than
report an eight-fold reduction against a quantity that has no interpretation. Its
Pearson $r$ is shown because the calibration is what makes the relationship linear.
\end{tablenotes}
\end{threeparttable}
\end{table}

\begin{figure}[H]
\centering
\includegraphics[width=\textwidth]{../../results/paper/figures/fig5_calibration.pdf}
\caption{\textbf{Monotone calibration restores an absolute efficiency scale.}
\textbf{(A)} The learned map $g$ from the ensemble's rank score to efficiency, fitted by
isotonic regression on out-of-fold development predictions only and then applied once
to the frozen held-out predictions.
\textbf{(B)} Calibration curve on held-out data, binned by predicted quantile; the
dotted line is perfect calibration.
\textbf{(C)} Metric comparison against the baseline on identical rows. Because $g$ is
non-decreasing it inverts no pair, so Spearman is unchanged; on the recovered
efficiency scale the model's absolute error is below the baseline's.}
\label{fig:calib}
\end{figure}

The residual $2\times10^{-5}$ change in Spearman arises because isotonic regression is
piecewise constant and therefore creates ties; no pair is inverted.

\subsection{The benchmark is not saturated}
\label{sec:headroom}

Using 649 technical replicate groups (\S\ref{sec:ceiling}) we estimate an empirical
repeatability ceiling for the Kim partition and compare it against the model on every
surface we score on (Table~\ref{tab:ceiling}). Averaged over the three matched
development folds the ceiling is $0.9061$ against a model score of $0.7912$, a gap of
$+0.1149$; on the held-out set it is $0.9167$ against $0.8124$, a gap of $+0.1043$.
Restricting to rows that actually edit leaves $+0.1046$ and $+0.0934$ respectively, so
the gap is not merely the zero block. The conclusion is the same on all five surfaces
and does not depend on which one is chosen: roughly a tenth of a Spearman unit of
reducible error remains on the benchmark's largest partition.

\begin{table}[H]
\centering
\begin{threeparttable}
\caption{Remaining headroom on the Kim partition against the replicate-based ceiling,
reported on every evaluation surface.}
\label{tab:ceiling}
\begin{tabular}{l S[table-format=5.0] S[table-format=1.4] S[table-format=1.4] S[table-format=+1.4]}
\toprule
Surface & {Rows} & {Model $\rhoS$} & {Ceiling $\rhoS$} & {Gap} \\
\midrule
\multicolumn{5}{l}{\emph{Matched development folds, out-of-fold}} \\
\quad fold 0                       & 19747 & 0.7869 & 0.9026 & +0.1157 \\
\quad fold 1                       & 19740 & 0.7883 & 0.9080 & +0.1198 \\
\quad fold 2                       & 19742 & 0.7985 & 0.9076 & +0.1092 \\
\quad \textbf{mean of three folds} & {} & \bfseries 0.7912 & \bfseries 0.9061 & {$\mathbf{+0.1149}$} \\
\quad mean, rows with $y>0$        & {} & 0.8472 & 0.9518 & +0.1046 \\
\midrule
\multicolumn{5}{l}{\emph{Official held-out set}} \\
\quad \textbf{frozen ensemble}     & 11334 & \bfseries 0.8124 & \bfseries 0.9167 & {$\mathbf{+0.1043}$} \\
\quad best single member           & 11334 & 0.8151 & 0.9165 & +0.1014 \\
\quad frozen ensemble, $y>0$       &  5928 & 0.8618 & 0.9552 & +0.0934 \\
\bottomrule
\end{tabular}
\begin{tablenotes}[flushleft]\footnotesize
\item Ceiling is $\sqrt{R}$, where $R$ is the replicate--replicate reliability
estimated from 649 groups in which all sixteen design and condition covariates and the
target site are identical (\S\ref{sec:ceiling}). The ceiling depends only on the target
distribution of the row set and not on any model, which is why it differs between
surfaces.
\item Replicates are resampled non-parametrically from 1{,}298 ordered replicate pairs,
which inherits zero-discordance, heteroscedasticity and boundedness from the data
rather than assuming a noise model.
\item Development-fold model scores are out-of-fold; held-out scores are the same
prediction vectors reported in Table~\ref{tab:main}. An earlier version of this
analysis reported the fold-0 row alone and labelled it ``all Kim rows''; we give all
three folds and both surfaces because no other claim in this paper rests on a single
fold (\S\ref{sec:resolution}).
\item Rows with $y>0$ show the gap is not an artifact of the zero block. The zero rows
alone carry no ordering to score, the target being constant on them.
\end{tablenotes}
\end{threeparttable}
\end{table}

Per-condition analysis (Figure~\ref{fig:ceiling}), shown for development fold 0 so that
each condition retains enough rows to estimate its own ceiling, finds the gap present in
every condition and scaling with zero-inflation: conditions with 60--64\% exact zeros
show gaps near $0.18$, while those near 38\% show gaps near $0.09$.

\begin{figure}[H]
\centering
\includegraphics[width=\textwidth]{../../results/paper/figures/fig6_ceiling.pdf}
\caption{\textbf{The benchmark is not saturated.}
\textbf{(A)} Model performance against the empirical repeatability ceiling for each Kim
condition on development fold 0, ordered by zero-inflation; the annotation gives the
gap. Headroom is present in every condition. The pooled gap is reported on all five
surfaces in Table~\ref{tab:ceiling}; this panel is per-condition and so is shown for a
single fold, where each condition retains enough rows to estimate a ceiling.
\textbf{(B)} The gap to ceiling increases with the fraction of rows measured as exactly
zero, the regime in which a single measurement is least informative.
\textbf{(C)} Both corrections described in \S\ref{sec:ceiling} are load-bearing. A
naive replicate key merges epegRNA with plain-pegRNA constructs and returns a ceiling
essentially equal to the model's own score (dashed line), which would wrongly indicate
saturation; treating observed zeros as noiseless instead over-estimates it. Only the
corrected, censoring-aware estimate is trustworthy.}
\label{fig:ceiling}
\end{figure}

\subsection{A tie-structure caveat for Spearman on zero-inflated data}
\label{sec:ties}

$28.4\%$ of matched development-fold rows have efficiency \emph{exactly} zero
($50.7\%$ within Kim, $0.7\%$ within Liu); the corresponding held-out figure is
$26.8\%$. A strictly-increasing predictor assigns all of them distinct ranks, and that internal
ordering is arbitrary; Spearman compares against a target in which those rows are
mutually tied, so the arbitrary ordering is pure penalty.

Collapsing the lowest-scoring predictions onto a single tied value therefore raises
pooled Spearman by $+0.0044$, positive on all development folds, with the optimal
fraction ($0.30$) coinciding with the development-fold zero-mass ($28.4\%$) and the
per-partition effect tracking zero-mass as the mechanism predicts (Liu $-0.0154$ at
$0.7\%$ zeros; Kim $+0.0066$ at $50.7\%$).

\textbf{We did not adopt this.} Applying the identical transform to five other models
yields $+0.0127$ to $+0.0162$ irrespective of model quality, so it exploits Spearman's
tie convention rather than improving prediction, and comparator models emit continuous
tie-free scores. We report it because it is a general hazard: \emph{on zero-inflated
targets, comparisons between predictors with different tie behaviour are not
automatically like-for-like}, and reported Spearman values should be interpreted with
the predictor's tie structure in mind.

\subsection{Interventions that did not help}
\label{sec:domain}

For completeness and to guide future work we summarise interventions that were tested
under matched controls and did not improve performance (Table~\ref{tab:neg}). Each was
evaluated against a mechanism-free control run in the same batch, and effects below
$\approx0.005$ were required to replicate across three development folds.

\begin{table}[H]
\centering
\begin{threeparttable}
\caption{Interventions tested under matched controls and not adopted.}
\label{tab:neg}
\begin{tabular}{l l S[table-format=+1.4] l}
\toprule
Category & Intervention & {$\Delta\rhoS$} & Reason not adopted \\
\midrule
\multirow{3}{*}{Post-hoc combination}
 & context-gated ensemble weights      & +0.0000 & null \\
 & nonlinear stacking                  & +0.0013 & below threshold \\
 & residual learning on features       & +0.0006 & oracle upper bound \\
\midrule
\multirow{4}{*}{Supervision}
 & auxiliary simplex head              & +0.0008 & matches control \\
 & multi-resolution ordinal            & +0.0008 & matches control \\
 & context-relative ordinal (primary)  & -0.0235 & harmful \\
 & quantile regression head            & -0.0140 & harmful \\
\midrule
\multirow{3}{*}{Parameterisation}
 & rank-consistent CORAL head          & +0.0005 & null \\
 & monotonicity penalty                & +0.0007 & null \\
 & hurdle (zero-inflation) head        & -0.0022 & harmful \\
\midrule
\multirow{5}{*}{Architecture}
 & selective (Mamba-style) SSM         & +0.0031 & below threshold \\
 & hybrid S4D$+$attention, 1:1 ratio   & -0.0027 & harmful \\
 & S4D state dimension 128             & -0.0022 & harmful \\
 & wider FFN / deeper stack / MoE      & {n.r.}  & did not replicate \\
 & layerwise context conditioning      & +0.0001 & null \\
\midrule
\multirow{3}{*}{Domain shift}
 & source loss re-weighting (mild)     & {n.r.}  & did not replicate \\
 & per-source output head              & -0.0018 & harmful \\
 & training on evaluated sources only  & -0.0294 & harmful \\
\bottomrule
\end{tabular}
\begin{tablenotes}[flushleft]\footnotesize
\item $\Delta\rhoS$ is measured against a matched control run in the same batch. For
capacity and parameterisation changes the control holds parameter count fixed; for the
per-source head it ties the heads; for supervision changes it is an unmodified re-run.
\item ``n.r.'' denotes a candidate that reached $+0.003$ to $+0.004$ on one development
fold but changed sign on another. Effects below $\approx0.005$ were required to
replicate across three folds before promotion, and none did.
\item Two entries are informative beyond their sign and are discussed in
\S\ref{sec:domain}: removing the unevaluated source study is strongly harmful, and
additional context conditioning increases rather than decreases context invariance.
\item The selective SSM makes $\Delta$, $B$ and $C$ functions of the input, which is
the mechanism separating Mamba \citep{gu2023mamba} from S4. Its entry is measured
against its own \emph{frozen-selection} control --- identical parameter count and
recipe with the selection projections held at initialisation, so the extra parameters
exist but cannot become content-dependent --- which is the comparison that isolates
selectivity. Both selective runs sit well below the S4D baseline they would replace
($0.8827$ and $0.8796$ against $0.8982$), but they sit below it \emph{together}, so
that deficit is attributable to the scan recipe rather than to selectivity: a selective
SSM cannot use the FFT convolution and needs a sequential scan, which forced batch 256,
fp32 and $\approx$12$\times$ the training time. The conclusion we draw is narrow --- a
selective mixer did not pay for itself here --- not that selectivity is useless on this
problem.
\end{tablenotes}
\end{threeparttable}
\end{table}

Two results in this table are informative beyond their sign.

First, \textbf{removing the unevaluated source study is strongly harmful}
($-0.0294$). The Schwank-derived partition contributes 58.4\% of training rows and 0\%
of the evaluation set, so discarding it appears to be free efficiency; it is not. Its
volume carries transferable signal that outweighs its distributional mismatch, and
aggressive re-weighting toward the evaluated distribution was likewise harmful.

Second, \textbf{additional context conditioning increases context invariance}. The
model's cross-condition rank correlation exceeds the empirical one ($0.858$ vs.\ a
disattenuated $0.758$), indicating it under-models how experimental context
\emph{reorders} designs. Applying FiLM at every block --- intended to address this ---
moved the diagnostic the wrong way ($0.828\to0.875$) and left performance unchanged.
FiLM is structurally a per-channel scale and shift, so additional conditioning capacity
improves representation of the cross-condition \emph{mean}, which explains far more
variance and is easier to fit, rather than creating reordering capacity.

% =====================================================================
\section{Discussion}

\subsection{Can architecture substitute for mechanism?}

For the specific task of ranking efficiency on this benchmark, yes. A model containing
no prime-editing reaction structure --- only a paired edit representation, a
segment-aware pegRNA representation, cross-attention between them, and categorical
context conditioning --- outperforms a strongly mechanistic predictor by $+0.0389$
Spearman pooled and $+0.0506$ within condition, under identical data and evaluation, and
also on absolute error (\S\ref{sec:calib-results}).

Three qualifications keep that from becoming a larger claim than the evidence supports.
First, the comparison covers one of the comparator's axes and none of the others
(\S\ref{sec:domain-weight} and the scope statement in \S1): accuracy on single-pegRNA
efficiency is not generalisation to PE3 or twinPE, and it is not interpretability. A
mechanistic parameterisation is a hypothesis about the reaction that can be inspected,
falsified and transferred; ours is not, and for a practitioner deciding \emph{why} a
design fails that difference may outweigh $0.04$ Spearman. Second, a substantial part
of the partition-specific margin is attributable to a training-weight choice rather
than to architecture (\S\ref{sec:domain-weight}). Third, the wider Kim margin
($+0.0804$) is on the partition spanning the most cell lines and editing systems, where
a prior tuned to a narrower regime would be expected to generalise least well; we
regard that as consistent with the reading rather than as demonstrating it, and the
weighting analysis supplies a more parsimonious explanation.

The one structural assumption we \emph{did} import proved useful, and is worth
distinguishing from reaction mechanism: the simplex head encodes that an edited locus
ends in exactly one of three measured outcomes. That is a property of the assay, not of
the biochemistry, and it supplies the indel signal a scalar head discards.

\subsection{Metric-matched supervision}

The ordinal head's contribution is best understood as aligning the training objective
with the evaluation. Predicting cumulative threshold indicators estimates a row's
normalised rank, which is precisely what a rank correlation scores, and it distributes
supervision evenly across a heavily skewed target instead of concentrating gradient on
a few high-efficiency rows.

The more transferable observation is that it changed \emph{which} rows the model gets
wrong, and this was predicted before the experiment rather than observed afterwards.
Residual correlation with the existing ensemble is about $0.05$ lower for ordinal-head
models than for models differing only in seed, capacity or data subsample, though the
two groups are not disjoint (\S\ref{sec:decorrelation}). For applications that ensemble
models, varying the objective is therefore a more reliable source of complementary
errors than varying capacity, seeds or data subsampling --- all of which we tested and
found ineffective. We would not, on this evidence, use residual correlation to screen an
individual candidate; we would use it to choose which axis to vary.

\subsection{State-space mixing for biological sequences}

That the state-space mixer contributes twice what the objective does deserves comment.
Our sequences are short (100 and 90 tokens), so the usual motivation for state-space
models --- linear-time long-context modelling \citep{gu2022s4,gu2023mamba} --- does not
apply. The benefit must therefore be representational. A diagonal state-space kernel
imposes smooth, distance-parameterised, translation-equivariant mixing, which matches a
domain where interactions are governed by relative position along the molecule
(PBS length, RTT length, distance from nick to edit) rather than by content-based
lookup among distant tokens.

This suggests state-space mixers deserve broader evaluation on short biological
sequences, where they are currently uncommon, and where their inductive bias may be a
better match than attention's. We note one negative constraint from our results: naive
hybridisation at a 1:1 ratio of attention to state-space blocks was harmful, consistent
with reports that successful hybrid language models use attention sparingly
\citep{lieber2024jamba}. Sparse hybridisation at published ratios remains untested here.

\subsection{What the remaining error is, and is not}

Our replicate analysis indicates roughly $+0.10$ Spearman of genuine headroom on the
Kim partition, consistently across all five surfaces we measured it on. That number required two corrections, each of which alone would have
reversed the conclusion --- excluding construct pairs masquerading as replicates, and
treating observed zeros as censored rather than certain. We highlight this because a
naive replicate analysis returns $0.7987$, essentially our model's score, and would
have supported a confident and incorrect claim of saturation.

Given the headroom is real, the natural question is why roughly seventy subsequent
experiments failed to capture it. Our diagnostic points at design $\times$ context
interaction: the model applies a nearly universal design ranking, while the data show
substantial context-dependent reordering. Direct attacks on this failed, including one
in which a mechanism-free shuffled-target control outperformed the genuine objective.
We think the most plausible remaining explanation is that the relevant covariates are
absent from the input: what distinguishes one cell line from another for a given design
--- mismatch-repair status, chromatin accessibility, repair-pathway balance --- is
supplied to the model only as a categorical label, and no objective over a categorical
label can recover biology the label does not encode. Testing this requires richer
context annotation rather than further architectural search.

\subsection{Limitations}

\begin{enumerate}
\item The absolute $0.90$ threshold is not established at 95\% confidence
($[0.8955,0.9173]$; 90\% of resamples above). The margin over the comparator is the
robust claim.
\item The ensemble does not outperform its best single member on held-out data; the
development-fold advantage did not transfer.
\item 196 of 649 replicate groups straddle the official train/held-out split, so some
held-out rows have an exact design-and-condition twin in training. This affects any
model trained on this split equally and so does not distort the comparison, but it
makes absolute held-out scores slightly optimistic.
\item The Kim partition has no unused evaluation surface remaining, so Kim-specific
claims rest on out-of-fold development folds and a once-consulted lockbox.
\item Our reproduction of the comparator scores above its published figure under every
aggregation convention we tried, and no evaluation script was released. The direction
favours the baseline and so does not undermine our claims, but it caps the precision of
the comparison.
\item\label{sec:limitations-baselines} All results are on a single benchmark, and
against one external comparator plus the feature baseline of \S\ref{sec:gbm}. External
validation on an independent prime-editing dataset is the most valuable next step and is
not yet available. We attempted two further baselines, PRIDICT2.0 and DeepPrime-FT, and
report neither: both require reproducing hand-engineered feature pipelines exactly ---
including melting temperatures, ViennaRNA folding energies \citep{lorenz2011} and a
separate pretrained on-target cutting model \citep{kim2019} --- and a subtly mismatched
reimplementation would produce misleading numbers rather than no numbers.
\item The comparison covers one of the comparator's capabilities. It does not address
PE3 or twinPE prediction, prospective design, or interpretability, all of which the
comparator provides and our model does not. A mechanistic parameterisation yields
inspectable per-step rates; ours yields a score.
\item The training-weight asymmetry in \S\ref{sec:domain-weight} is inferred from the
released code applied to the files the authors supplied, not verified against their
training procedure directly. If the released checkpoints were in fact trained with
different weights, the attribution in that section changes even though the measured
comparison does not.
\item The held-out set was consulted four times across model generations rather than
once (\S\ref{sec:resolution}). Within-generation decisions were made on development
folds and a once-consulted lockbox, and the paired design protects the margin, but the
absolute level carries some adaptive-selection exposure.
\end{enumerate}

\subsection{Methodological recommendations}

Three practices materially changed our conclusions and we recommend them for work at
this effect size.

\textbf{Include a mechanism-free control in every experimental batch.} In one round,
fourteen candidates each gained $+0.0017$ to $+0.0020$; an unmodified re-run of the
same model at a different batch size gained $+0.0018$. The difference attributable to
every idea in that round was $+0.0002$. In another, a control with \emph{shuffled}
supervision targets outperformed the genuine objective.

\textbf{Require replication before believing sub-$0.005$ effects.} Three interventions
on orthogonal axes each reached $+0.003$ to $+0.004$ on one development fold and all
three changed sign on a second.

\textbf{State the mechanism as a falsifiable prediction.} Our context-conditioning
experiment was accompanied by a pre-registered prediction about the cross-condition
diagnostic. The prediction failed in a specific direction, which identified the reason
(FiLM strengthens the mean-shift pathway) far faster than the null performance result
alone would have.

% =====================================================================
\section{Conclusions}

A minimally mechanistic relational model matches and exceeds a strongly mechanistic
predictor of prime-editing efficiency on its own benchmark, reaching Spearman $0.9079$
against $0.8690$ ($\Delta=+0.0389$, ahead in all 5{,}000 protospacer-clustered bootstrap
resamples). A factorial analysis attributes roughly two-thirds of the gain to replacing
self-attention with a bidirectional state-space sequence mixer and one-third to an
ordinal outcome head matched to the rank-correlation metric and to the target's zero
inflation; a single model containing both matches a five-member ensemble. Monotone
calibration fitted purely on out-of-fold data restores an interpretable absolute
efficiency estimate without altering the ranking, and on that scale the model also
leads the baseline ($0.0478$ against $0.0590$ mean absolute error).

The benchmark is not saturated: replicate analysis indicates roughly $+0.10$ Spearman
of reducible error remains on its largest partition. Our evidence suggests the
obstacle is not model capacity or supervision --- roughly seventy controlled
experiments across architecture, objective, parameterisation and training scheme failed
to close it --- but the absence of covariates describing how experimental context
reorders designs. Progress will likely require richer context annotation rather than
further architectural search.

% =====================================================================
\section*{Data and code availability}

All code, configurations, per-run metadata, the complete research log including every
negative result, and the frozen model specification are maintained in the project
repository. The training corpus was provided by the authors of the comparator method
and is not redistributed here.

\section*{Acknowledgements}

We thank the authors of the comparator method for providing their exact training mix,
without which a matched comparison would not have been possible.

% =====================================================================
\begin{thebibliography}{99}\small

\bibitem[Anzalone et al.(2019)]{anzalone2019}
Anzalone AV, Randolph PB, Davis JR, et al.
Search-and-replace genome editing without double-strand breaks or donor DNA.
\emph{Nature} \textbf{576}, 149--157 (2019).

\bibitem[Cao et al.(2020)]{cao2020}
Cao W, Mirjalili V, Raschka S.
Rank consistent ordinal regression for neural networks with application to age
estimation. \emph{Pattern Recognition Letters} \textbf{140}, 325--331 (2020).

\bibitem[Chen et al.(2021)]{chen2021}
Chen PJ, Hussmann JA, Yan J, et al.
Enhanced prime editing systems by manipulating cellular determinants of editing
outcomes. \emph{Cell} \textbf{184}, 5635--5652 (2021).

\bibitem[Doench et al.(2016)]{doench2016}
Doench JG, Fusi N, Sullender M, et al.
Optimized sgRNA design to maximize activity and minimize off-target effects of
CRISPR-Cas9. \emph{Nature Biotechnology} \textbf{34}, 184--191 (2016).

\bibitem[Efron \& Tibshirani(1993)]{efron1993}
Efron B, Tibshirani RJ.
\emph{An Introduction to the Bootstrap.} Chapman \& Hall, New York (1993).

\bibitem[Gu et al.(2022a)]{gu2022s4}
Gu A, Goel K, R\'e C.
Efficiently modeling long sequences with structured state spaces.
\emph{International Conference on Learning Representations} (2022).

\bibitem[Gu et al.(2022b)]{gu2022s4d}
Gu A, Gupta A, Goel K, R\'e C.
On the parameterization and initialization of diagonal state space models.
\emph{Advances in Neural Information Processing Systems} (2022).

\bibitem[Gu \& Dao(2023)]{gu2023mamba}
Gu A, Dao T.
Mamba: Linear-time sequence modeling with selective state spaces.
\emph{arXiv:2312.00752} (2023).

\bibitem[Hsu et al.(2026)]{hsu2026}
Hsu A, Chen P, Li A, Hemez C, Gao X, Terrey M, Nelson C, Selvam V, Cristian A,
McElroy A, Steinbeck B, Mahadeshwar G, Pandey S, Barsdale Z, Chen P, Sousa A,
Sakai H, Silverstein R, Morad I, Krueger R, Shen M, Kleinstiver B, Lutz C,
Tolar J, Blazar B, Osborn M, Liu DR.
Mechanistic machine learning for prediction of prime editing outcomes.
\emph{Nature Biotechnology} \textbf{44} (2026).
\url{https://doi.org/10.1038/s41587-026-03261-7}.

\bibitem[Kim et al.(2019)]{kim2019}
Kim HK, Kim Y, Lee S, et al.
SpCas9 activity prediction by DeepSpCas9, a deep learning-based model with high
generalization performance. \emph{Science Advances} \textbf{5}, eaax9249 (2019).

\bibitem[Kim et al.(2021)]{kim2021}
Kim HK, Yu G, Park J, et al.
Predicting the efficiency of prime editing guide RNAs in human cells.
\emph{Nature Biotechnology} \textbf{39}, 198--206 (2021).

\bibitem[Lieber et al.(2024)]{lieber2024jamba}
Lieber O, Lenz B, Bata H, et al.
Jamba: a hybrid Transformer-Mamba language model. \emph{arXiv:2403.19887} (2024).

\bibitem[Lorenz et al.(2011)]{lorenz2011}
Lorenz R, Bernhart SH, H\"oner zu Siederdissen C, et al.
ViennaRNA Package 2.0. \emph{Algorithms for Molecular Biology} \textbf{6}, 26 (2011).

\bibitem[Mathis et al.(2023)]{mathis2023}
Mathis N, Allam A, Kissling L, et al.
Predicting prime editing efficiency and product purity by deep learning.
\emph{Nature Biotechnology} \textbf{41}, 1151--1159 (2023).

\bibitem[Mathis et al.(2024)]{mathis2024}
Mathis N, Allam A, Tálas A, et al.
Machine learning prediction of prime editing efficiency across diverse chromatin
contexts. \emph{Nature Biotechnology} (2024).

\bibitem[P\'erez et al.(2018)]{perez2018}
P\'erez E, Strub F, de Vries H, Dumoulin V, Courville A.
FiLM: visual reasoning with a general conditioning layer.
\emph{AAAI Conference on Artificial Intelligence} (2018).

\bibitem[Spearman(1904)]{spearman1904}
Spearman C.
The proof and measurement of association between two things.
\emph{American Journal of Psychology} \textbf{15}, 72--101 (1904).

\bibitem[Vaswani et al.(2017)]{vaswani2017}
Vaswani A, Shazeer N, Parmar N, et al.
Attention is all you need.
\emph{Advances in Neural Information Processing Systems} (2017).

\bibitem[Yu et al.(2023)]{yu2023}
Yu G, Kim HK, Park J, et al.
Prediction of efficiencies for diverse prime editing systems in multiple cell types.
\emph{Cell} \textbf{186}, 2256--2272 (2023).

\bibitem[Zadrozny \& Elkan(2002)]{zadrozny2002}
Zadrozny B, Elkan C.
Transforming classifier scores into accurate multiclass probability estimates.
\emph{ACM SIGKDD International Conference on Knowledge Discovery and Data Mining}
(2002).

\end{thebibliography}

\end{document}
